\chapter{Kurzweil--Henstock Integration}


% Section numbers and titles follow the board (original/9.11, photo 40).
% Board sections 1--4 correspond to book sections 2.1--2.4.
\section{Cells, divisions \& partial divisions}

\begin{definitionenv}
An interval $I\subseteq\mathbb R$ is \textbf{non-degenerate} if
$I\neq\varnothing$ and $I$ is not a singleton. A \textbf{cell} is a
non-degenerate interval.
\end{definitionenv}

\begin{notationenv}
For $A\subseteq\mathbb R$, write
\[
    \operatorname{int}(A),\qquad
    \operatorname{clos}(A),\qquad
    \operatorname{bdry}(A)
\]
for the interior, closure, and boundary of $A$.
\end{notationenv}

\begin{propositionenv}
A cell is compact if and only if it is closed and bounded.
\end{propositionenv}

\begin{definitionenv}
For nonempty $S\subseteq\mathbb R$,
\[
    \operatorname{diam}S
    :=
    \sup\{|x-y|:x,y\in S\}.
\]
\end{definitionenv}

\begin{definitionenv}
A sequence $(I_n)_{n\in\mathbb N}$ is a sequence of \textbf{nested cells} if
\[
    I_{n+1}\subseteq I_n\quad\forall n,
    \qquad
    \operatorname{diam}I_n\to0,
\]
and every $I_n$ is a compact cell.
\end{definitionenv}

\begin{theoremenv}
If $(I_n)_{n\in\mathbb N}$ is a sequence of nested cells, then
$\bigcap_{n\in\mathbb N}I_n$ is a singleton.
\end{theoremenv}

\begin{exerciseenv}
Let $J$ be a cell and $x\in J$. Show that there exists a compact cell
$J'\subseteq J$ such that $x\notin J'$.
\end{exerciseenv}

\begin{theoremenv}
Every cell is uncountable.
\end{theoremenv}

\begin{proofenv}
Suppose $I=\{x_n:n\in\mathbb N\}$.
Using the preceding exercise, construct nested compact cells $(I_n)$ such that
\[
    I_0\subseteq I,\qquad x_n\notin I_n,\qquad
    \operatorname{diam}I_{n+1}<\frac12\operatorname{diam}I_n.
\]
Then
\[
    \bigcap_{n\in\mathbb N}I_n=\{x\}\subseteq I.
\]
Since $I=\{x_n:n\in\mathbb N\}$, we have $x=x_m$ for some $m$, contradicting
$x_m\notin I_m$.
\end{proofenv}


\begin{definitionenv}
Two cells $I_1,I_2$ are \textbf{non-overlapping} if
\[
    \operatorname{clos}(I_1)\cap\operatorname{clos}(I_2)
\]
is either empty or a singleton.
\end{definitionenv}

\begin{propositionenv}
The following basic properties hold.

\begin{enumerate}[label=(\Alph*)]
    \item If $I$ is a cell, then
    \[
        \operatorname{int}(\operatorname{clos}I)
        \subseteq I
        \subseteq
        \operatorname{clos}(\operatorname{int}I).
    \]

    \item If $I_1,I_2$ are non-overlapping closed cells, then
    \[
        \operatorname{int}(I_1)\cap I_2=\varnothing.
    \]

    \item If $I_1,I_2$ are cells with $I_1\cap I_2\ne\varnothing$,
    then $I_1\cup I_2$ is a cell.

    \item Let $I_1,I_2$ be non-overlapping closed cells. The following are
    equivalent:
    \begin{enumerate}[label=(\arabic*)]
        \item $I_1\cup I_2$ is a cell;
        \item $I_1\cap I_2$ is a singleton;
        \item $I_1\cap I_2=\{x\}$ for some $x\in\operatorname{int}(I_1\cup I_2)$;
        \item $I_1\cap I_2\neq\varnothing$.
    \end{enumerate}

    \item If $I_1,I_2,I_3$ are pairwise non-overlapping, then
    \[
        I_1\cap I_2\cap I_3=\varnothing.
    \]

    \item If $I_1,I_2$ are disjoint cells, then
    $\operatorname{clos}(I_1)$ and $\operatorname{clos}(I_2)$ are non-overlapping.
\end{enumerate}
\end{propositionenv}

\begin{lemmaenv}
Let $I$ be a compact cell and $N\in\mathbb N_+$. Suppose
$I_1,\ldots,I_N$ are pairwise non-overlapping compact cells and
\[
    I=\bigcup_{n=1}^{N}I_n.
\]
After renumbering if necessary,
\[
\begin{gathered}
    \min I_1=\min I,\qquad \max I_N=\max I,\\
    \max I_n=\min I_{n+1}\qquad(1\le n<N).
\end{gathered}
\]
\end{lemmaenv}

\begin{exerciseenv}
Prove the preceding lemma.
\end{exerciseenv}

\begin{definitionenv}
For a compact cell $[a,b]$, define
\[
    \lambda([a,b])=b-a.
\]
\end{definitionenv}

\begin{definitionenv}
A \textbf{division} of a compact cell $I$ is a finite collection
$\Delta$ of pairwise non-overlapping compact cells such that
\[
    \bigcup\Delta=I.
\]
If $\Delta=\{I_1,\ldots,I_N\}$, then
\[
    \bigcup\Delta:=\bigcup_{n=1}^{N}I_n.
\]
\end{definitionenv}

\begin{propositionenv}
If $\Delta=\{I_1,\ldots,I_N\}$ is a division of $I$, then
\[
    \lambda(I)=\sum_{n=1}^{N}\lambda(I_n).
\]
\end{propositionenv}

\begin{proofenv}
Renumber the cells as in the preceding lemma and write $I_n=[a_n,b_n]$.
Then
\[
\begin{aligned}
    \sum_{n=1}^{N}\lambda(I_n)
    &=(b_1-a_1)+\cdots+(b_N-a_N)\\
    &=\max I-\min I\\
    &=\lambda(I).
\end{aligned}
\]
\end{proofenv}

\begin{theoremenv}[Borel]
Let $I,I_1,I_2,\ldots$ be compact cells. If
\[
    I\subseteq\bigcup_{n=1}^{\infty}I_n,
\]
then
\[
    \sum_{n=1}^{\infty}\lambda(I_n)\ge\lambda(I).
\]
\end{theoremenv}

\begin{definitionenv}
A \textbf{partial division} of a compact cell $I$ is a finite collection
$\Delta$ of pairwise non-overlapping compact cells such that
\[
    \bigcup\Delta\subseteq I.
\]
\end{definitionenv}

\begin{notationenv}
\[
    \mathcal C(I)
    :=
    \{J:J\text{ is a compact cell contained in }I\}.
\]
\end{notationenv}


\Needspace{10\baselineskip}
\section{Gauges, fine P-divisions \& Cousin theorem}

\begin{definitionenv}
A \textbf{gauge} on a compact cell $I$ is a function
$\delta:I\to\mathbb R_{>0}$.
For gauges $\delta_1,\ldots,\delta_N$ on $I$, their pointwise minimum
\[
    \delta:=\min\{\delta_1,\ldots,\delta_N\},\qquad
    \delta(x)=\min_{1\le k\le N}\delta_k(x)
\]
is also a gauge on $I$.
\end{definitionenv}

\Needspace{7\baselineskip}
\begin{definitionenv}
A \textbf{$P$-division} of $I$ (respectively, a
\textbf{partial $P$-division}) is a finite set
\[
    \Pi\subseteq\mathcal C(I)\times I
\]
such that
\begin{enumerate}[label=(\arabic*)]
    \item if $(J,x)\in\Pi$, then $x\in J$;
    \item if $(J_1,x_1),(J_2,x_2)\in\Pi$ are distinct, then
    $J_1,J_2$ are non-overlapping;
    \item
    \[
        \Delta_\Pi
        :=
        \{J:(J,x)\in\Pi\text{ for some }x\}
    \]
    is a division (respectively, a partial division) of $I$.
\end{enumerate}
\end{definitionenv}

\begin{definitionenv}
Let $\delta$ be a gauge on $I$. A (partial) $P$-division $\Pi$ is
\textbf{$\delta$-fine} if
\[
    J\subseteq[x-\delta(x),x+\delta(x)]
\]
for every $(J,x)\in\Pi$.
\end{definitionenv}

% Lecture 2026-09-11: original/9.11, photographs 40--43.
\begin{notationenv}
The \textbf{body} of a partial $P$-division $\Pi$ is
\[
    \operatorname{body}(\Pi)
    :=\bigcup\{J:(J,x)\in\Pi\}=\bigcup\Delta_\Pi.
\]
\end{notationenv}

\begin{propositionenv}
Let $I$ be a compact cell.
\begin{enumerate}[label=(\Alph*)]
    \item Suppose $\Pi_1,\ldots,\Pi_N$ are partial $P$-divisions of $I$
    and, for $j\ne k$, the intersection
    $\operatorname{body}(\Pi_j)\cap\operatorname{body}(\Pi_k)$
    is empty or a singleton. Then $\bigcup_{k=1}^N\Pi_k$ is a partial
    $P$-division of $I$.
    \item For any partial $P$-division $\Pi$ and $\xi\in I$, at most two
    pairs $(J,x)\in\Pi$ have $x=\xi$.
\end{enumerate}
\end{propositionenv}


\Needspace{9\baselineskip}
\begin{exampleenv}
If a gauge $\delta$ on $I=[a,b]$ is bounded below by $d>0$, choose
$N\in\mathbb N_+$ with $(b-a)/N<d$. Put
\[
    a_n=a+\frac{n(b-a)}N\quad(0\le n\le N),
    \qquad J_n=[a_n,a_{n+1}]\quad(0\le n<N).
\]
Choose any $x_n\in J_n$. Then
$\Pi=\{(J_n,x_n):0\le n<N\}$ is a $\delta$-fine $P$-division:
for $y\in J_n$, $|y-x_n|\le(b-a)/N<\delta(x_n)$.
\end{exampleenv}

\begin{lemmaenv}\label{lem:gluing-tagged-divisions}
Let $\Delta$ be a division of a compact cell $I$. For each $K\in\Delta$,
let $\Pi_K$ be a $P$-division of $K$. Then
\[
    \bigcup_{K\in\Delta}\Pi_K
\]
is a $P$-division of $I$. If each $\Pi_K$ is $\delta|_K$-fine for a gauge
$\delta$ on $I$, their union is $\delta$-fine.
\end{lemmaenv}

\begin{exerciseenv}
Prove Lemma~\ref{lem:gluing-tagged-divisions}.
\end{exerciseenv}

\begin{theoremenv}[Cousin]\label{thm:cousin}
Given a compact cell $I$ and a gauge $\delta$ on $I$, there exists a
$\delta$-fine $P$-division of $I$.
\end{theoremenv}

\begin{proofenv}
Suppose no such division exists. Bisect $I_0=I$ into two compact cells.
At least one half has no $\delta$-fine $P$-division; otherwise the gluing
lemma would give one for $I_0$. Choose such a half as $I_1$ and repeat.
This gives nested compact cells satisfying
\[
    I_{n+1}\subseteq I_n,\qquad
    \operatorname{diam}I_n=2^{-n}\operatorname{diam}I,
\]
none of which admits a $\delta|_{I_n}$-fine $P$-division.
The nested-cell theorem gives
\[
    \bigcap_{n=0}^{\infty}I_n=\{x\}.
\]
Since $\delta(x)>0$, choose $n_0$ such that
$\operatorname{diam}I_{n_0}<\delta(x)$. As $x\in I_{n_0}$,
\[
    I_{n_0}\subseteq[x-\delta(x),x+\delta(x)].
\]
Thus $\{(I_{n_0},x)\}$ is a $\delta|_{I_{n_0}}$-fine $P$-division,
a contradiction.
\end{proofenv}


% Board section 3, title transcribed from photo 42.
\section{KH-int., lin., mon. \& Cauchy crit.}
% Lecture 2026-09-11: photographs 42--48.

\begin{notationenv}
For a compact cell $I$, we also write its length as
\[
    |I|:=\lambda(I)=\max I-\min I.
\]
Thus $|[a,b]|=b-a$, and $J\subseteq I$ implies $|J|\le|I|$.
For a division $\Delta$ of $I$, the length identity proved earlier reads
\[
    \sum_{J\in\Delta}|J|=|I|.
\]
\end{notationenv}

\begin{definitionenv}
Let $f:I\to\mathbb R$ and let $\Pi$ be a partial $P$-division of $I$.
The \textbf{Riemann sum} of $f$ over $\Pi$ is
\[
    S(\Pi,f):=\sum_{(J,x)\in\Pi}f(x)|J|.
\]
In particular, $S(\varnothing,f)=0$.
\end{definitionenv}

\begin{propositionenv}
Riemann sums have the following properties.
\begin{enumerate}[label=(\Alph*)]
    \item For $\alpha\in\mathbb R$,
    \[
        S(\Pi,\alpha f_1+f_2)=\alpha S(\Pi,f_1)+S(\Pi,f_2).
    \]
    \item If $f_1\le f_2$ pointwise on $I$, then
    \[
        S(\Pi,f_1)\le S(\Pi,f_2).
    \]
    \item If $\operatorname{body}(\Pi_1)\cap\operatorname{body}(\Pi_2)$
    is empty or a singleton, then
    \[
        S(\Pi_1\cup\Pi_2,f)=S(\Pi_1,f)+S(\Pi_2,f).
    \]
\end{enumerate}
\end{propositionenv}

\begin{proofenv}
\textbf{(C).} If $(J,x)\in\Pi_1\cap\Pi_2$, then
$J\subseteq\operatorname{body}(\Pi_1)\cap\operatorname{body}(\Pi_2)$,
contradicting non-degeneracy. Hence $\Pi_1\cap\Pi_2=\varnothing$, and
\[
    S(\Pi_1\cup\Pi_2,f)
    =\sum_{(J,x)\in\Pi_1}f(x)|J|
     +\sum_{(J,x)\in\Pi_2}f(x)|J|.
\]
\end{proofenv}

\Needspace{9\baselineskip}
\begin{definitionenv}
A function $f:I\to\mathbb R$ is \textbf{Kurzweil--Henstock integrable}
(\textbf{KH-integrable}) on the compact cell $I$ if there exists
$r\in\mathbb R$ such that, for every $\varepsilon>0$, there is a gauge
$\delta$ on $I$ for which
\[
    |S(\Pi,f)-r|<\varepsilon
\]
for every $\delta$-fine $P$-division $\Pi$ of $I$.
\end{definitionenv}

\begin{lemmaenv}
The number $r$ in the definition of KH-integrability is unique.
\end{lemmaenv}

\begin{proofenv}
Let $r_1,r_2$ satisfy the definition and fix $\varepsilon>0$.
Choose corresponding gauges $\delta_1,\delta_2$, and set
$\delta=\min\{\delta_1,\delta_2\}$. By Cousin's theorem, there is a
$\delta$-fine $P$-division $\Pi$ of $I$. It is fine for both gauges, so
\[
    |r_1-r_2|
    \le |r_1-S(\Pi,f)|+|S(\Pi,f)-r_2|<2\varepsilon.
\]
Since $\varepsilon$ is arbitrary, $r_1=r_2$.
\end{proofenv}

\begin{notationenv}
The unique number above is the \textbf{KH integral} of $f$ on $I$, denoted by
\[
    \int_I f=\int_I f(x)\,\dd x.
\]
For $I=[a,b]$, we also write $\int_a^b f$ or $\int_a^b f(x)\,\dd x$.
The variable of integration is a dummy variable; for instance,
\[
    \int_I f(x)\,\dd x=\int_I f(u)\,\dd u.
\]
We write
\[
    \mathrm{KH}(I):=\{f:I\to\mathbb R:f\text{ is KH-integrable on }I\}.
\]
A gauge $\delta$ is \textbf{$\varepsilon$-adapted to $f$} if every
$\delta$-fine $P$-division $\Pi$ satisfies
$|S(\Pi,f)-\int_I f|<\varepsilon$.
\end{notationenv}

\begin{exampleenv}
If $f(x)=c$ for all $x\in I$, then for every $P$-division $\Pi$,
\[
    S(\Pi,f)=c\sum_{(J,x)\in\Pi}|J|=c|I|.
\]
Thus $f\in\mathrm{KH}(I)$ and $\int_I f=c|I|$.
\end{exampleenv}

\begin{remark}
KH-integrable functions on compact intervals include constant functions,
step functions, continuous functions, monotone functions,
Riemann-integrable functions, Lebesgue-integrable functions, and derivatives.
\end{remark}

\Needspace{12\baselineskip}
\begin{theoremenv}[Linearity and monotonicity]\label{thm:kh-linearity-monotonicity}
The set $\mathrm{KH}(I)$ is a real vector space, and the map
\[
    \mathrm{KH}(I)\longrightarrow\mathbb R,
    \qquad f\longmapsto\int_I f
\]
is linear. Moreover, if $f_1,f_2\in\mathrm{KH}(I)$ and $f_1\le f_2$, then
\[
    \int_I f_1\le\int_I f_2.
\]
\end{theoremenv}

\begin{proofenv}
\textbf{Monotonicity.} For $\varepsilon>0$, choose
$\varepsilon$-adapted gauges $\delta_1,\delta_2$ and a
$\min\{\delta_1,\delta_2\}$-fine $P$-division $\Pi$ by Cousin's theorem.
Since $S(\Pi,f_1)\le S(\Pi,f_2)$,
\[
\begin{aligned}
    \int_I f_1-\int_I f_2
    &\le \int_I f_1-S(\Pi,f_1)+S(\Pi,f_2)-\int_I f_2\\
    &\le\left|\int_I f_1-S(\Pi,f_1)\right|
       +\left|S(\Pi,f_2)-\int_I f_2\right|<2\varepsilon.
\end{aligned}
\]
Letting $\varepsilon\downarrow0$ gives the assertion.
\end{proofenv}


% The Cauchy criterion and subcell corollary remain in board section 3.
% Lecture 2026-09-11: photographs 47--49.

\begin{theoremenv}[Cauchy criterion]\label{thm:kh-cauchy}
Let $I$ be a compact cell and $f:I\to\mathbb R$. The following are equivalent:
\begin{enumerate}[label=(\Alph*)]
    \item $f\in\mathrm{KH}(I)$;
    \item for every $\varepsilon>0$, there is a gauge $\delta$ on $I$ such
    that, for any two $\delta$-fine $P$-divisions $\Pi,\widetilde\Pi$ of $I$,
    \[
        |S(\Pi,f)-S(\widetilde\Pi,f)|<\varepsilon.
    \]
\end{enumerate}
\end{theoremenv}

\begin{proofenv}
\textbf{$(A)\Rightarrow(B)$.} Choose an $\varepsilon/2$-adapted gauge.
For any two divisions fine for this gauge,
\[
    |S(\Pi,f)-S(\widetilde\Pi,f)|
    \le\left|S(\Pi,f)-\int_I f\right|
       +\left|\int_I f-S(\widetilde\Pi,f)\right|<\varepsilon.
\]

\textbf{$(B)\Rightarrow(A)$.} For each $n\in\mathbb N_+$, choose a gauge
$\delta_n$ for the error $1/n$ in (B), and define
\[
    \widehat\delta_n:=\min\{\delta_1,\ldots,\delta_n\}.
\]
By Cousin's theorem, choose a $\widehat\delta_n$-fine $P$-division $\Pi_n$
of $I$, and put $r_n=S(\Pi_n,f)$. If $p<q$, both $\Pi_p$ and $\Pi_q$
are $\delta_p$-fine, so
\[
    |r_p-r_q|<\frac1p.
\]
Thus $(r_n)$ is Cauchy in $\mathbb R$ and converges to some $r\in\mathbb R$.

Given $\varepsilon>0$, choose $n_0$ such that $1/n_0<\varepsilon/2$.
Let $\Pi$ be any $\delta_{n_0}$-fine $P$-division of $I$.
For $n\ge n_0$, $\Pi_n$ is also $\delta_{n_0}$-fine. Hence
\[
\begin{aligned}
    |r-S(\Pi,f)|
    &\le |r-r_n|+|S(\Pi_n,f)-S(\Pi,f)|\\
    &<|r-r_n|+\frac1{n_0}.
\end{aligned}
\]
Letting $n\to\infty$ gives
$|r-S(\Pi,f)|\le1/n_0<\varepsilon$.
Therefore $f\in\mathrm{KH}(I)$ and $\int_I f=r$.
\end{proofenv}

\begin{corollaryenv}[Subcells]\label{cor:kh-subcells}
Let $I,J$ be compact cells with $J\subseteq I$. If $f\in\mathrm{KH}(I)$,
then $f|_J\in\mathrm{KH}(J)$.
\end{corollaryenv}

\begin{proofenv}
Write $I=[a,b]$ and $J=[c,d]$. If $a<c$, set $J_1=[a,c]$;
if $d<b$, set $J_2=[d,b]$. Omit any degenerate end interval.

Fix $\varepsilon>0$ and choose a gauge $\delta$ on $I$ that is
$\varepsilon/2$-adapted to $f$. By Cousin's theorem, fix
$\delta|_{J_k}$-fine $P$-divisions $\Pi_k$ on the end intervals present.
Let $\Pi,\widetilde\Pi$ be any two $\delta|_J$-fine $P$-divisions of $J$.
Adjoin the same fixed end divisions to obtain $\Pi_0,\widetilde\Pi_0$,
both $\delta$-fine $P$-divisions of $I$.
Additivity of Riemann sums cancels the end contributions, giving
\[
\begin{aligned}
    |S(\Pi,f|_J)-S(\widetilde\Pi,f|_J)|
    &=|S(\Pi_0,f)-S(\widetilde\Pi_0,f)|\\
    &\le\left|S(\Pi_0,f)-\int_I f\right|
       +\left|\int_I f-S(\widetilde\Pi_0,f)\right|\\
    &<\varepsilon.
\end{aligned}
\]
The Cauchy criterion gives $f|_J\in\mathrm{KH}(J)$.
\end{proofenv}


% Board section 4, title transcribed from photo 50.
\section{Functions that agree except on a countable set}
% Lecture 2026-09-11: photographs 49--50 and 52--53.

\begin{definitionenv}
The \textbf{Dirichlet function} is
\[
    \mathbf1_{\mathbb Q}:\mathbb R\to\mathbb R,\qquad
    \mathbf1_{\mathbb Q}(x)=
    \begin{cases}
        1,&x\in\mathbb Q,\\
        0,&x\notin\mathbb Q.
    \end{cases}
\]
\end{definitionenv}

\begin{propositionenv}
For every compact cell $I$, we have
$\mathbf1_{\mathbb Q}|_I\in\mathrm{KH}(I)$ and
\[
    \int_I\mathbf1_{\mathbb Q}=0.
\]
\end{propositionenv}

\begin{proofenv}
Fix $\varepsilon>0$ and choose a bijection
$q:\mathbb N_+\to\mathbb Q\cap I$. Define a gauge on $I$ by
\[
    \delta(x)=
    \begin{cases}
        7^{58},&x\notin\mathbb Q,\\
        \displaystyle\frac{\varepsilon}{2^{n+2}},&x=q(n).
    \end{cases}
\]
Let $\Pi$ be a $\delta$-fine $P$-division and let
$F=\{n:q(n)\text{ occurs as a tag in }\Pi\}$, a finite set.
Only rational tags contribute. Each tag occurs at most twice, and
each associated cell satisfies $|J|\le2\delta(x)$. Consequently,
\[
\begin{aligned}
    0\le S(\Pi,\mathbf1_{\mathbb Q}|_I)
    &=\sum_{\substack{(J,x)\in\Pi\\x\in\mathbb Q}}|J|\\
    &\le4\sum_{n\in F}\delta(q(n))
      =4\sum_{n\in F}\frac{\varepsilon}{2^{n+2}}\\
    &<4\sum_{n=1}^{\infty}\frac{\varepsilon}{2^{n+2}}
      =\varepsilon.
\end{aligned}
\]
Therefore $\int_I\mathbf1_{\mathbb Q}=0$.
\end{proofenv}

\begin{theoremenv}\label{thm:kh-countable-changes}
Let $I$ be a compact cell and $f,g:I\to\mathbb R$. If
$f\in\mathrm{KH}(I)$ and
\[
    \{f\ne g\}:=\{x\in I:f(x)\ne g(x)\}
\]
is countable, then $g\in\mathrm{KH}(I)$ and
\[
    \int_I g=\int_I f.
\]
\end{theoremenv}
